\documentclass[11pt]{article}

\usepackage[margin=1in]{geometry}
\usepackage{amsmath}
\usepackage{booktabs}
\usepackage{enumitem}
\usepackage[hidelinks]{hyperref}

\newcommand{\code}[1]{\texttt{#1}}

\title{Wasted Work in Gabriel: Push, Sections, or Gated Chains}
\author{Design summary}
\date{\today}

\begin{document}
\maketitle

\begin{abstract}
Gabriel today is a \emph{push} design: every engine works on the newest
input it can get and passes its result on, whether or not the next engine
can use it. This keeps results fresh, but wastes work: an engine can compute
a result that a slower engine after it then skips. This note compares push
with two ways to avoid that waste: per-section token pools (``sections'')
and gated chains. The recommendation is to keep push as the default, and to
add gated chains as an option: an engine starts an input only when the
slowest engine after it will be ready for its result. In simulation of
nine chains, gated chains wasted no work, kept results about as fresh as
push, and needed no token estimates. They cost some throughput, more with
jittery engines. Sections also wasted no work, but were never better than
gated chains, and were clearly worse on chains with several slow engines.
All numbers are from simulation and bench runs on one machine; fan-outs and
remote crossings are not simulated yet.
\end{abstract}

\section{Where we are}

\begin{description}[style=nextline, leftmargin=1.5em, nosep]
  \item[Built] The client computes each producer's token count from the
        flow. A spec no longer needs one, and a count given too low gets a
        warning. The bench has a shadow estimator and a simulator
        (\code{bench/shadow\_counts.py}, \code{bench/simulate\_sections.py}),
        and run files for the sweeps.
  \item[Decided] Do not build sections (\S\ref{sec:sections}).
  \item[Proposed] Keep push as the default. Add gated chains
        (\S\ref{sec:gated}) as an option per flow. Simulated on nine
        chains; not built.
  \item[Open] Simulations of the paper's own pipelines (\S\ref{sec:paper}),
        fan-outs, a device engine feeding a crossing, filters inside a
        gated chain, how early a stretch should start under jitter, and a
        bench with a shared accelerator (\S\ref{sec:next}).
\end{description}

\section{Why this matters for SteelEagle}

In SteelEagle (\code{ddil2026}), the DroneHub runs on-drone AI under tight
size, weight, and power limits, and usually has one AI accelerator (GPU, NPU,
or TPU) shared by every on-drone engine. Work that is computed and then
thrown away costs battery and accelerator time that another engine could
have used. On the cloudlet, compute is plentiful, but the link to it is
DDIL: its bandwidth and round trip change, and it can drop.

The paper describes Gabriel's flow control as ``token-based'', ``ensuring
that delays cause data to be dropped rather than queued'', and adapting to
the bottleneck ``regardless of whether it is in transmission, processing, or
multi-tenant queuing.'' Any change has to keep those properties. This rules
out designs that queue inputs, which matters for sections below.

\section{Today's design: push}
\label{sec:push}

Gabriel today pushes: no engine waits for the engine after it to be ready.

\begin{itemize}[nosep]
  \item \textbf{Tokens.} Each producer has a pool of tokens. An input takes
        one when it enters the flow and gives it back at its first
        completion. The client computes the count from the flow: one token
        per engine and per edge on the producer's paths, the capacity of
        each crossing to the cloudlet, and one spare (``Choosing $C$'' in
        \code{dag-tokens.tex}).
  \item \textbf{Newest wins.} Each edge has one pending slot. A newer input
        replaces the one waiting there. So an engine always starts on the
        newest input it can get.
  \item \textbf{Crossings} to the cloudlet have a capacity: how many inputs
        may be on the link or waiting at the cloudlet. This is the window
        for the network.
\end{itemize}
Push keeps results fresh, needs no estimates, does not care if there are
more tokens than needed, and absorbs jitter: when an engine runs long, the
next one still finds a recent input waiting. Its one cost is wasted work.
That cost is real only when the wasted work takes something another engine
needs: a shared accelerator, a battery, or a cloudlet GPU shared by several
drones. When every engine has its own hardware, push is the right design.

\section{Where wasted work comes from}

Work is wasted when an engine computes a result that nobody uses. That
happens in only one place: \textbf{an engine in a chain whose results go
only to a slower engine}, directly or through faster ones. In a chain
P $\to$ A $\to$ B with A at 50\,ms and B at 500\,ms, A produces ten results
for each one B can take, and B skips the other nine. In simulation of that
chain, 45\% of all engine time went to results B skipped. A slow engine can
itself be wasteful, if a still slower one comes after it.

Other places look similar but waste nothing:
\begin{itemize}[nosep]
  \item \textbf{The producer.} A camera frame that is replaced before any
        engine starts on it costs nothing.
  \item \textbf{A fan-out.} If F feeds a fast engine C and a slow branch,
        C uses every one of F's results. The slow branch skips results
        before any work on them.
  \item \textbf{An engine the app reads.} If the app subscribes to A's
        results, A's work is used even when B skips it. In SteelEagle, the
        on-drone filter is such an engine: its scores drive thresholding,
        and its data is logged and used during disconnected operation.
\end{itemize}
So the problem is local to chains. A fix only needs to act inside chains.

\section{Options considered}

\subsection{A lower token count}

With one token, A only works on inputs B will take. In the chain above, that
saved 90\% of A's work at 7.5\% of B's throughput. But it fails elsewhere:
\begin{itemize}[nosep]
  \item In a chain of three equal engines, one token gives a third of the
        throughput.
  \item Two tokens can freeze a slot: an input waits in B's slot holding the
        last token, and nothing newer can replace it. B then starts on
        inputs 525\,ms old.
  \item At a fan-out, the token comes back when the fastest branch finishes,
        so a slow branch is fed just as fast with one token as with ten.
\end{itemize}
One count for the whole flow cannot fit every shape.

\subsection{Sections}
\label{sec:sections}

\paragraph{The idea.} Apply the original Gabriel again at every fan-out. A
\emph{section} is a chain from a producer or a fan-out up to the next
fan-out. Each section has its own token pool. An input enters a section only
with one of its tokens, and returns it at its first completion in the
section; this rule is recursive over nested fan-outs. Inside a section,
inputs wait in order, so nothing is skipped after work was spent on it.

\paragraph{The count.} Each pool would set its count while the flow runs,
from the knee of a closed queueing network:
\[
  N = \max\left(1,\ \left\lfloor \frac{D + Z}{\max(D_{\max}, I)}
  + \epsilon \right\rfloor\right)
\]
where $D$ is the total service time per input (each engine's time weighted by
the share of inputs that reach it), $Z$ the network round trips, $D_{\max}$
the slowest engine's share, $I$ the producer's interval, and $\epsilon$ a
small tolerance. The client can measure all of these from events it already
gets. A shadow run on recorded bench runs gave stable estimates, the same
whatever token count the run used.

\paragraph{Why not sections.}
\begin{enumerate}[nosep]
  \item \textbf{They queue.} Inputs that wait in order age while they wait.
        That goes against ``dropped rather than queued''.
  \item \textbf{A count is the wrong control.} Freshness depends on
        \emph{when} an engine starts, and a count only sets \emph{how many}
        inputs are in a chain. The right count is often a fraction, so it
        must be rounded: down leaves an engine idle, up makes inputs wait.
        With two slow engines in a chain, the fraction is 1.48: one token
        loses a third of the throughput, and two make results about 150\,ms
        older (Table~\ref{tab:several}).
  \item \textbf{A count that is too high is costly.} In order, each token
        above the right count adds one whole time of the slowest engine to
        every input's wait: 500\,ms per token in the chain above. Shared
        engines, a heavy tail, or a changing round trip all push the
        estimate the wrong way. With newest wins, extra tokens cost nothing.
  \item \textbf{Their fan-out machinery solves no waste.} Branch pools and
        nested returns act at fan-outs, which waste nothing.
  \item \textbf{Much more code}: pools on the client and on servers, queues
        on edges, a central recount, a control call to push counts to
        servers.
\end{enumerate}
Sections keep one real advantage: every input that enters a section is
processed by all of it, which suits strict fusion. Gabriel has no joins yet.

\subsection{Gated chains}
\label{sec:gated}

\paragraph{The rule.} Keep push, and add one condition. For an
engine in a chain, let S be the slowest engine after it in the chain. If S
is slower than this engine:
\begin{quote}
Start an input only if the next slot is empty, and no earlier than the time
S will be ready for a new input, minus the service times of the engines in
between.
\end{quote}
If no engine after it is slower, the engine pushes, as today: nothing after it
can skip its results.
\begin{itemize}[nosep]
  \item \textbf{The gate} (next slot empty) means a result is never computed
        that will be skipped: in a chain, nothing else can fill that slot.
        This alone guarantees no wasted work.
  \item \textbf{The timing} keeps results fresh. Without it, A would start
        as soon as B's slot empties, and its result would wait in B's slot
        for all of B's time.
  \item \textbf{Inputs on their way to S count.} If B already carries the
        input S takes next, A's result is for the run after that. A first
        version that looked only one engine ahead got this wrong.
  \item \textbf{Time to the slowest engine after, not the last one.} With S
        in the middle of a chain, timing to the last engine made results
        1308\,ms old instead of 852\,ms.
\end{itemize}
Each engine needs only its median service time, measured on its own server,
and the state of the engines between it and S.

\paragraph{Example.} P $\to$ A $\to$ B $\to$ S with A at 20\,ms, B at 30\,ms,
and S at 500\,ms. S starts input $x$ at time 0 and will be free at 500. B
starts at $500 - 30 = 470$, and A at $470 - 20 = 450$. Until then, A's slot
keeps the newest camera frame. S starts on a frame about 50\,ms old, and A
and B run once per run of S.

\paragraph{Several slow engines.} Walk the chain from its end, keeping the
slowest engine seen so far. An engine slower than everything after it is an
\emph{anchor}, and pushes, as today. Every other engine is timed to the slowest
engine after it. So a chain splits into stretches, each ending at an anchor:
\begin{center}
\small
\begin{tabular}{@{}ll@{}}
\toprule
Chain (ms) & Timing \\
\midrule
A 20 $\to$ S1 200 $\to$ B 20 $\to$ S2 500 & A, S1, B to S2 \\
A 20 $\to$ S1 500 $\to$ B 20 $\to$ S2 200 & A to S1; S1 free; B to S2 \\
A 20 $\to$ S1 100 $\to$ B 20 $\to$ S2 300 $\to$ C 20 $\to$ S3 200 &
  A, S1, B to S2; S2 free; C to S3 \\
\bottomrule
\end{tabular}
\end{center}
When two slow engines are equal, noise in their measured medians picks the
anchor. Either choice works.

\paragraph{Where a chain ends.} A gated chain stops at a fan-out, a join, an
engine the app reads, or a crossing to another server. Beyond those points,
push applies. So the rule never has to reason about the whole graph:
more branches only mean more short chains.

\paragraph{The cloudlet.} Timing a send to arrive just when a cloudlet
engine frees up would need a prediction across a DDIL link, and a wrong
guess leaves the cloudlet's GPU idle. A window does better there: the
crossing capacity, sized to the round trip (the original Gabriel's tokens).
The earlier simulator in \code{sim/} agrees: at 150\,ms delay, the modes
with tokens on the link (\code{token}, \code{token-adaptive}, and
\code{segment}) gave an AoI of 705--734\,ms, and the timed ones
(\code{schedule}, \code{schedule-hop}) 780--868\,ms. So:
\begin{itemize}[nosep]
  \item a chain inside the cloudlet is gated on the cloudlet, as on the
        device;
  \item the link keeps its window;
  \item a device engine whose only consumer is a crossing is gated on the
        crossing. The crossing's ``free time'' is when its oldest input is
        expected back: its send time, plus the round trip, plus the cloudlet
        engine's service time. The client knows all three on its own clock.
\end{itemize}
This last case is not simulated yet.

\section{Results}

The simulator draws service times from those measured in bench runs, and
matched the measured runs within 6\,ms of AoI and 2\% of throughput
(\code{bench/simulate\_sections.py} for push and sections; the gated
rule is in a scratch copy for now). Each cell gives the mean AoI at the last
engine in ms, results a second, and the share of engine time wasted; where
no waste is shown, it was zero. Push and the gated chains use the
computed token count; sections use their own count, as the rule gives it.

\begin{table}[h]
\centering
\small
\begin{tabular}{@{}llll@{}}
\toprule
Chain (ms) & Push & Sections & Gated chain \\
\midrule
One heavy-tailed engine  & 118 & 118 & 118 \\
Three equal engines      & 76  & 76  & 76 \\
Bench chain (filter, detector, tracker) & 105 & 111 & 110--112 \\
A 50 $\to$ B 500         & 843; 2.0; 45\% & 843; 1.8 & 817; 2.0 \\
A 20 $\to$ B 30 $\to$ S 500 & 829; 2.0; 56\% & 836; 1.8 & 844; 1.95 \\
S 500 $\to$ A 20 $\to$ B 30 & 820; 2.0 & --- & 810; 2.0 \\
\bottomrule
\end{tabular}
\caption{Chains with one slow engine.}
\label{tab:one}
\end{table}

\begin{table}[h]
\centering
\small
\begin{tabular}{@{}llll@{}}
\toprule
Chain (ms), 10\% jitter & Push & Sections & Gated chain \\
\midrule
A 20, S 500, C 20, D 30 (S in the middle) & 835; 2.01; 33\% & 871; 1.76 &
  852; 1.98 \\
A 20, S1 200, B 20, S2 500 (rising) & 1108; 2.01; 45\% & 1116; 1.36 &
  \textbf{1035}; 1.92 \\
A 20, S1 500, B 20, S2 200 (falling) & 1014; 1.99; 28\% & 1122; 1.36 &
  1032; 1.97 \\
A 20, S1 500, B 20, S2 500 (equal) & 1435; 1.97; 23\% & 1364; 1.86 &
  1379; 1.89 \\
Three slow engines (100, 300, 200) & 878; 3.34; 38\% & 857; 2.98 &
  \textbf{846}; 3.18 \\
\midrule
\multicolumn{4}{@{}l}{With 30\% jitter:} \\
Rising  & 1153; 1.98; 45\% & 1110; 1.40 & 1088; 1.80 \\
Falling & 1046; 1.98; 28\% & 1120; 1.37 & 1088; 1.95 \\
Equal   & 1508; 1.78; 25\% & 1545; 1.66 & 1541; 1.71; 1\% \\
Three   & 918; 3.26; 38\%  & 918; 2.81  & 908; 2.94 \\
\bottomrule
\end{tabular}
\caption{Chains with several slow engines. Sections use the count their
rule gives (1 for rising, falling, and the middle case; 2 for equal and
three). With one more token, sections reach push's throughput, but their
AoI rises by 140--400\,ms.}
\label{tab:several}
\end{table}

\begin{itemize}[nosep]
  \item \textbf{Gated chains waste nothing} (at most 1\%), against up to
        56\% with push.
  \item \textbf{Their freshness is about push's}: from 7\% fresher to 7\%
        older. They are fresher when a slow engine is itself wasteful, as in
        the rising chain.
  \item \textbf{They cost throughput}: 1--5\% with 10\% jitter, up to 10\%
        with 30\%. Timing a stretch adds up the service times of every
        engine in it, so their errors add up too. When a stretch runs long,
        its slow engine waits.
  \item \textbf{Their AoI does not depend on the token count}, so the
        computed count can stay.
  \item \textbf{Push keeps the most throughput}, and handles jitter best.
        Its waste is 23--56\% of engine time in every chain with a slow
        engine after a faster one, and none otherwise.
  \item \textbf{Sections were never better than gated chains} on both
        measures at once. With several slow engines, their count falls
        between two whole numbers, and either choice loses: a third of the
        throughput, or 140--400\,ms of freshness.
\end{itemize}

\section{Comparison}

\begin{center}
\small
\begin{tabular}{@{}p{3.2cm}p{3.5cm}p{3.5cm}p{3.5cm}@{}}
\toprule
 & Push (today) & Sections & Gated chains \\
\midrule
Wasted work & 23--56\% in chains with a slow engine after a faster one &
  None & None \\
AoI & Baseline & 5\% fresher to 11\% older & 7\% fresher to 7\% older \\
Throughput & Baseline & Up to a third lower & 1--10\% lower, more with
  jitter \\
Jitter & Absorbed by waiting inputs & Costs throughput or freshness &
  Costs throughput \\
Inputs queue? & No & Yes & No \\
Estimates & None & Service times, visit shares, round trips, interval &
  Median service time per engine \\
A wrong estimate costs & --- & A whole bottleneck time per extra token &
  About the size of the error \\
Too many tokens & Free & Costly & Free \\
Decisions made & Locally & Centrally, pushed to servers & Locally \\
New code & --- & Pools, queues, recount, control call &
  A start condition in the server's pull loop \\
\bottomrule
\end{tabular}
\end{center}

\paragraph{Which to use.}
\begin{itemize}[nosep]
  \item \textbf{Push} when engines have their own hardware, or the waste is
        cheap. It is the default.
  \item \textbf{Gated chains} when wasted work takes a shared resource: a
        shared accelerator, a battery, a shared cloudlet GPU. Set per flow.
  \item \textbf{Sections} never: gated chains do as well or better, without
        queues or count estimates.
\end{itemize}

\section{For the paper}
\label{sec:paper}

\paragraph{The paper's pipelines.} The simulations above use made-up
chains, not SteelEagle's pipelines. Mapped to Gabriel flows, the paper's
modes are:
\begin{center}
\small
\begin{tabular}{@{}p{3.3cm}p{6.3cm}p{4.3cm}@{}}
\toprule
Mode & Flow & Wasted work under push \\
\midrule
Strongly connected ($>$5\,Mbps) & camera $\to$ cloudlet detector &
  None: no engines on the drone \\
Weakly connected (30\,kbps--5\,Mbps) & camera (lower rate) $\to$ tiling and
  on-drone filter $\to$ crossing $\to$ cloudlet detector &
  Not in the filter, whose results are logged and used. Only in a separate
  tiler or encoder before it \\
Multi-tier fusion & camera $\to$ \{on-drone detector, cloudlet detector\}
  $\to$ app & None: a fan-out. The on-drone engines may share one
  accelerator \\
Disconnected ($<$30\,kbps) & camera $\to$ on-drone AI &
  Only if that AI is a chain that ends in a slower engine \\
\bottomrule
\end{tabular}
\end{center}
So, as the paper describes SteelEagle, push wastes little on the drone, and
gated chains are a refinement, not a core mechanism. They matter if the
pipeline grows chains: a tiler or encoder before the filter, a detector
before a slow classifier, or a cloudlet chain shared by a drone swarm.

\paragraph{What matters more for the paper,} and is not simulated yet:
\begin{itemize}[nosep]
  \item \textbf{The crossing under DDIL}: bandwidth from 30\,kbps to
        5\,Mbps, mode switches with 20\% hysteresis, a changing round
        trip, and outages. How the crossing capacity should adapt.
  \item \textbf{A shared on-drone accelerator}, used by the filter and the
        fusion detector at once. Gabriel does not model or schedule it.
  \item \textbf{Tiling}: one 4K frame becomes many 256$\times$256 tiles, so
        the filter is a one-to-many engine. Gabriel's \code{Emit} sends one
        output per input.
\end{itemize}
The \code{sim/} package already has an emulated link with outages,
fan-outs, filters, and joins, so it is the place for these scenarios. It
needs measured timings: on-drone and cloudlet model latencies, frame and
tile sizes, and the cellular round trip.

\paragraph{What the paper can claim.} The paper's description of Gabriel
stays true: token-based, drop rather than queue, adapts to the bottleneck.
In addition:
\begin{itemize}[nosep]
  \item \textbf{The token count is computed from the flow}, so a pipeline
        author does not tune it, and a wrong count cannot freeze a slot.
  \item \textbf{The link keeps a window} sized to the round trip, as in the
        original Gabriel.
  \item \textbf{Gated chains} (once built and measured) avoid wasted work
        in on-drone chains without making results older.
\end{itemize}
The comparison of push, sections, and gated chains could be a short design
discussion. A claim about gated chains needs a bench with a shared
accelerator, so that the saved work shows up as fresher or more frequent
results.

\section{Next steps}
\label{sec:next}

\begin{enumerate}[nosep]
  \item Simulate the paper's pipelines in \code{sim/}, with measured
        timings, comparing push with gated chains where they apply.
  \item Move the gated rule into \code{bench/simulate\_sections.py}, so these
        tables can be reproduced from the repository.
  \item How early a stretch should start under jitter. Starting from a
        pessimistic sum of service times (say the 80th percentile) instead
        of the medians should recover throughput, at some cost in freshness.
  \item Filters inside a gated chain. ``Slowest'' should compare each
        engine's time weighted by the share of inputs that reach it. A
        filter that passes inputs at random makes the timing less exact:
        it cost 5\% of AoI in the bench chain.
  \item Simulate a fan-out with a chain branch (F $\to$ \{C, A $\to$ B\}):
        only A is gated, and C should stay as fresh as with push.
  \item Simulate a device engine gated on a crossing, with a changing round
        trip and an outage.
  \item Add a shared accelerator to the bench, so wasted work has a cost.
  \item If these hold, build the gate in the server's pull loop, behind a
        setting per flow, and measure it on the bench and over the real
        link.
  \item Update \code{dag-tokens.tex}: replace the Sections subsection with
        gated chains, and keep sections as a rejected alternative.
\end{enumerate}

\end{document}
